\section{引言：为什么需要序列建模}

本讲是 MIT 6.S191（Introduction to Deep Learning）的第二讲，由 Ava Amini 主讲。在上一讲中，Alexander Amini 介绍了神经网络的基础知识，包括感知机（Perceptron）、前馈网络（Feedforward Network）以及反向传播（Backpropagation）算法。本讲将从这些基础出发，系统地讲解如何将神经网络应用于\textbf{序列数据}（Sequential Data）的处理，涵盖三个核心主题：

\begin{enumerate}
    \item \textbf{Recurrent Neural Networks (RNNs)}：循环神经网络的基本原理与内部机制
    \item \textbf{训练RNN的挑战}：梯度消失/爆炸问题及 LSTM 等改进架构
    \item \textbf{Attention 机制与 Transformer}：自注意力机制的原理及其在现代大语言模型中的应用
\end{enumerate}

\begin{importantbox}{序列建模的核心思想}
序列建模的核心目标是：给定一个序列的历史信息，预测该序列未来的状态。这是大语言模型（LLM）、语音识别、机器翻译等众多 AI 应用的基础。
\end{importantbox}

\subsection{直觉引入：预测球的运动}

讲者用一个生动的例子引入序列建模的概念：假设在二维空间中有一个球，我们的任务是预测它接下来的运动方向。

\begin{figure}[H]
\centering
\includegraphics[width=0.7\textwidth]{slides-images/slide-03.jpg}
\caption{仅凭球的当前位置无法确定运动方向\protect\footnotemark}
\end{figure}
\footnotetext{Slides 第3页。}

如果只给出球当前的位置而没有任何历史信息，我们的预测只能是随机的——球可能向任何方向移动。但是，如果给出球过去几个时间步的位置序列，问题就变得简单多了：

\begin{figure}[H]
\centering
\includegraphics[width=0.7\textwidth]{slides-images/slide-05.jpg}
\caption{利用历史位置序列可以准确预测运动方向\protect\footnotemark}
\end{figure}
\footnotetext{Slides 第5页。}

通过观察球从左到右的历史轨迹，我们可以合理推断它将继续向右移动。这个例子揭示了序列建模的核心：\textbf{利用过去的信息来预测未来}。

\subsection{序列数据无处不在}

序列数据在现实世界中无处不在：

\begin{figure}[H]
\centering
\includegraphics[width=0.85\textwidth]{slides-images/slide-07.jpg}
\caption{现实世界中的序列数据类型\protect\footnotemark}
\end{figure}
\footnotetext{Slides 第7页。}

\begin{itemize}
    \item \textbf{音频信号}：语音波形可以被分割成一系列声波片段进行序列处理
    \item \textbf{自然语言}：文本可以被拆分为单词或字符的序列
    \item \textbf{医学信号}：如心电图（ECG）的时间序列
    \item \textbf{金融数据}：股票价格的时间序列
    \item \textbf{生物序列}：DNA/RNA 核酸序列、蛋白质氨基酸序列
    \item \textbf{天气数据、视频帧序列}等
\end{itemize}

\subsection{序列建模的任务类型}

\begin{knowledgebox}{序列建模的三种基本模式}
\begin{itemize}
    \item \textbf{Many-to-One}：序列输入 $\rightarrow$ 单一输出。例如：情感分类（输入一段文本，输出正面/负面）
    \item \textbf{One-to-Many}：单一输入 $\rightarrow$ 序列输出。例如：图像描述生成（输入一张图片，输出描述文字）
    \item \textbf{Many-to-Many}：序列输入 $\rightarrow$ 序列输出。例如：机器翻译（英文序列 $\rightarrow$ 中文序列）
\end{itemize}
\end{knowledgebox}

\begin{figure}[H]
\centering
\includegraphics[width=0.75\textwidth]{slides-images/slide-08.jpg}
\caption{序列建模的不同任务类型\protect\footnotemark}
\end{figure}
\footnotetext{Slides 第8页。}

\subsection{本章小结}

序列建模是深度学习中最重要的范式之一，其核心在于利用数据的时间/顺序依赖关系。从简单的球运动预测到复杂的语言模型，序列建模的思想贯穿了现代 AI 的方方面面。

\newpage
